\documentclass{beamer}

\pdfmapfile{+sansmathaccent.map}


\mode<presentation>
{
	\usetheme{Warsaw} % or try Darmstadt, Madrid, Warsaw, Rochester, CambridgeUS, ...
	\usecolortheme{seahorse} % or try seahorse, beaver, crane, wolverine, ...
	\usefonttheme{serif}  % or try serif, structurebold, ...
	\setbeamertemplate{navigation symbols}{}
	\setbeamertemplate{caption}[numbered]
} 


%%%%%%%%%%%%%%%%%%%%%%%%%%%%
% itemize settings


%%%%%%%%%%%%%%%%%%%%%%%%%%%%
% itemize settings

\definecolor{myhotpink}{RGB}{255, 80, 200}
\definecolor{mywarmpink}{RGB}{255, 60, 160}
\definecolor{mylightpink}{RGB}{255, 80, 200}
\definecolor{mypink}{RGB}{255, 30, 80}
\definecolor{mydarkpink}{RGB}{155, 25, 60}

\definecolor{mypaleblue}{RGB}{240, 240, 255}
\definecolor{mylightblue}{RGB}{120, 150, 255}
\definecolor{myblue}{RGB}{90, 90, 255}
\definecolor{mygblue}{RGB}{70, 110, 240}
\definecolor{mycourseblue}{RGB}{120, 140, 255}
\definecolor{mycourselightblue}{RGB}{245, 250, 255}
\definecolor{mydarkblue}{RGB}{0, 0, 180}
\definecolor{myblackblue}{RGB}{40, 40, 120}

\definecolor{mygreen}{RGB}{0, 200, 0}
\definecolor{mygreen2}{RGB}{245, 255, 230}
\definecolor{mydarkgreen}{RGB}{0, 120, 0}
\definecolor{mydarkgreen}{RGB}{0, 120, 0}


\definecolor{mygray}{gray}{0.8}
\definecolor{mydarkgray}{RGB}{80, 80, 160}
\definecolor{mylightgray}{RGB}{160, 160, 160}

\definecolor{mydarkred}{RGB}{160, 30, 30}
\definecolor{mylightred}{RGB}{255, 150, 150}
\definecolor{myred}{RGB}{200, 110, 110}
\definecolor{myblackred}{RGB}{120, 40, 40}

\definecolor{mypink}{RGB}{255, 30, 80}
\definecolor{myhotpink}{RGB}{255, 80, 200}
\definecolor{mywarmpink}{RGB}{255, 60, 160}
\definecolor{mylightpink}{RGB}{255, 80, 200}
\definecolor{mydarkpink}{RGB}{155, 25, 60}
\definecolor{mywhitepink}{RGB}{255, 240, 240}

\definecolor{mydarkcolor}{RGB}{60, 25, 155}
\definecolor{mylightcolor}{RGB}{130, 180, 250}



\setbeamertemplate{itemize items}[default]

\setbeamertemplate{itemize item}{\color{myblackblue}$\blacksquare$}
\setbeamertemplate{itemize subitem}{\color{mygblue}$\blacktriangleright$}
\setbeamertemplate{itemize subsubitem}{\color{mygray}$\blacksquare$}

\setbeamercolor{palette quaternary}{fg=white,bg=mycourseblue}
\setbeamercolor{titlelike}{parent=palette quaternary}

\setbeamercolor{palette quaternary2}{fg=black,bg=mycourselightblue}
\setbeamercolor{frametitle}{parent=palette quaternary2}

\setbeamerfont{frametitle}{size=\Large,series=\scshape}
\setbeamerfont{framesubtitle}{size=\normalsize,series=\upshape}





%%%%%%%%%%%%%%%%%%%%%%%%%%%%
% block settings

\setbeamercolor{block title}{bg=red!30,fg=black}

\setbeamercolor*{block title example}{bg=mygreen!40!white,fg=black}

\setbeamercolor*{block body example}{fg= black, bg= mygreen2}


%%%%%%%%%%%%%%%%%%%%%%%%%%%%
% URL settings
\hypersetup{
	colorlinks=true,
	linkcolor=blue,
	filecolor=blue,      
	urlcolor=blue,
}

%%%%%%%%%%%%%%%%%%%%%%%%%%

\renewcommand{\familydefault}{\rmdefault}


\newcommand{\mycolor}[1] {
\setbeamercolor{normal text}{fg={#1}} % Recolors standard body text
\setbeamercolor{structure}{fg={#1}}   % Recolors bullets and structure items
\usebeamercolor[fg]{normal text}      % Applies the normal text color immediately
}


\usepackage{amsmath}
\usepackage{mathtools}

\usepackage{cancel}
\usepackage{centernot}

\usepackage{changepage}
\usepackage{subcaption}

\usepackage{qrcode}

\newcommand{\bo}[1] {\mathbf{#1}}
\newcommand{\R}{\mathbb{R}} 
\newcommand{\T}{^\top}     



\newcommand{\mydate}{Fall 2026}

\newcommand{\mygit}{\textcolor{blue}{\href{https://github.com/SergeiSa/Advanced-Control-2026}{github.com/SergeiSa/Advanced-Control-2026}}}

\newcommand{\myqr}{ \textcolor{black}{\qrcode[height=1.5in]{https://github.com/SergeiSa/Advanced-Control-2026}}
}

\newcommand{\myqrframe}{
	\begin{frame}
		\centerline{Lecture slides are available via Github:}
		\bigskip
		\centerline{\mygit}
		\bigskip
		\myqr
	\end{frame}
}


\newcommand{\bref}[2] {\textcolor{blue}{\href{#1}{#2}}}

%%%%%%%%%%%%%%%%%%%%%%%%%%%%
% code settings

\usepackage{listings}
\usepackage{color}
% \definecolor{mygreen}{rgb}{0,0.6,0}
% \definecolor{mygray}{rgb}{0.5,0.5,0.5}
\definecolor{mymauve}{rgb}{0.58,0,0.82}
\lstset{ 
	backgroundcolor=\color{white},   % choose the background color; you must add \usepackage{color} or \usepackage{xcolor}; should come as last argument
	basicstyle=\footnotesize,        % the size of the fonts that are used for the code
	breakatwhitespace=false,         % sets if automatic breaks should only happen at whitespace
	breaklines=true,                 % sets automatic line breaking
	captionpos=b,                    % sets the caption-position to bottom
	commentstyle=\color{mygreen},    % comment style
	deletekeywords={...},            % if you want to delete keywords from the given language
	escapeinside={\%*}{*)},          % if you want to add LaTeX within your code
	extendedchars=true,              % lets you use non-ASCII characters; for 8-bits encodings only, does not work with UTF-8
	firstnumber=0000,                % start line enumeration with line 0000
	frame=single,	                   % adds a frame around the code
	keepspaces=true,                 % keeps spaces in text, useful for keeping indentation of code (possibly needs columns=flexible)
	keywordstyle=\color{blue},       % keyword style
	language=Octave,                 % the language of the code
	morekeywords={*,...},            % if you want to add more keywords to the set
	numbers=left,                    % where to put the line-numbers; possible values are (none, left, right)
	numbersep=5pt,                   % how far the line-numbers are from the code
	numberstyle=\tiny\color{mygray}, % the style that is used for the line-numbers
	rulecolor=\color{black},         % if not set, the frame-color may be changed on line-breaks within not-black text (e.g. comments (green here))
	showspaces=false,                % show spaces everywhere adding particular underscores; it overrides 'showstringspaces'
	showstringspaces=false,          % underline spaces within strings only
	showtabs=false,                  % show tabs within strings adding particular underscores
	stepnumber=2,                    % the step between two line-numbers. If it's 1, each line will be numbered
	stringstyle=\color{mymauve},     % string literal style
	tabsize=2,	                   % sets default tabsize to 2 spaces
	title=\lstname                   % show the filename of files included with \lstinputlisting; also try caption instead of title
}


%%%%%%%%%%%%%%%%%%%%%%%%%%%%
% URL settings
\hypersetup{
	colorlinks=false,
	linkcolor=blue,
	filecolor=blue,      
	urlcolor=blue,
}

%%%%%%%%%%%%%%%%%%%%%%%%%%

%%%%%%%%%%%%%%%%%%%%%%%%%%%%
% tikz settings

\usepackage{tikz}
\tikzset{every picture/.style={line width=0.75pt}}
\usetikzlibrary{shapes.geometric, decorations.markings, positioning, patterns, calc, arrows.meta}


\subtitle{Advanced Control}
\author{by Sergei Savin}
\centering
\date{\mydate}



\title{Adjoint Method, Sensitivity Propagation}
\begin{document}
\maketitle


\begin{frame}{Content}
	\begin{itemize}
		\item Path planning on graphs
		\item Dijkstra's algorithm: initialization, exploring and updating nodes, selecting the next node, recovering the path
		\item A$^\ast$ algorithm and admissible heuristics
		\item Value iteration and the Bellman equation
	\end{itemize}
\end{frame}


%------------------------------------------------
\begin{frame}{Path Planning on Graphs}
	\begin{flushleft}
		
		Many practical path planning problems can be
		formalized using graphs:
		
		\begin{itemize}
			\item We describe the environment as a set of
			locations (nodes).
			\item We know between which locations we can
			move (edges).
			\item We know the cost of moving from one
			location to another (edge cost).
		\end{itemize}
		
		Thus, path planning becomes a problem of finding
		a sequence of edges connecting the starting node
		to the goal node, such that the sum of the edge
		costs is as small as possible.
		
	\end{flushleft}
\end{frame}

%------------------------------------------------
\begin{frame}{Paths on a Graph}
	\begin{flushleft}
		
		For any node $k$, let $\mathcal{N}(k)$ denote
		the set of nodes that can be reached from $k$
		by traversing one edge. These are called
		\textbf{adjacent nodes}.
		
		\medskip
		
		A path can be represented as a sequence of nodes
		connected by edges. Each intermediate node has
		a predecessor and a successor.
		
		
		\medskip
		
		
		An optimal path has an important property:
		
		\begin{itemize}
			\item Each intermediate node is reached from
			its predecessor along an optimal path to that
			node.
			\item From each intermediate node, the remaining
			path to the goal is also optimal.
		\end{itemize}
		
		
	\end{flushleft}
\end{frame}

%------------------------------------------------
\begin{frame}{Dijkstra's Algorithm}
	\begin{flushleft}
		
		\textbf{Dijkstra's algorithm} finds shortest paths
		from a starting node to other nodes in a graph.
		It can also be used to find a shortest path to a
		specific goal node.
		
		
		\medskip
		The key idea is to maintain the best known
		cost of reaching each discovered node from
		the starting node.
		
		\medskip
		
		We maintain two sets:
		
		\begin{itemize}
			\item \textbf{Open set:} nodes that have been
			discovered but not yet fully explored.
			\item \textbf{Closed set:} nodes whose shortest
			path costs have been finalized.
		\end{itemize}
		
		We assume that all edge costs are nonnegative.
		
	\end{flushleft}
\end{frame}

%------------------------------------------------
\begin{frame}{Dijkstra's Algorithm: Initialization}
	\begin{flushleft}
		
		Initially, only the starting node $s$ has been
		discovered:
		
		\[
		\text{Open}=\{s\},
		\qquad
		\text{Closed}=\emptyset.
		\]
		
		For every discovered node $k$, we maintain a
		tentative cost $g(k)$: the cost of the best path
		to $k$ found so far. The initial cost is
		\[
		g(s)=0.
		\]
		
		All other nodes are initially undiscovered, with
		their costs treated as infinity:
		\[
		g(k)=+\infty,\qquad k\ne s.
		\]
		
		We also store the predecessor of each discovered
		node to reconstruct the path.
		
	\end{flushleft}
\end{frame}

%------------------------------------------------
\begin{frame}{Dijkstra's Algorithm: Exploring a Node}
	\begin{flushleft}
		
		Exploring a node $k$ involves the following steps:
		
		\begin{enumerate}
			\item Check whether $k$ is the goal node.
			\item Examine every adjacent node $m\in\mathcal{N}(k)$.
			\item Move $k$ from the open set to the
			closed set.
		\end{enumerate}
		
		
		\medskip
		When examining an adjacent node $m$, we check
		whether the path through $k$ improves its
		best known cost. If $m$ has not been discovered, its tentative
		cost is initialized as
		\[
		g(m)=g(k)+c(k,m),
		\]
		where $c(k,m)$ is the cost of the edge from
		$k$ to $m$.
		
		We also record $k$ as the predecessor of $m$
		and add $m$ to the open set.
		
	\end{flushleft}
\end{frame}

%------------------------------------------------
\begin{frame}{Dijkstra's Algorithm: Updating a Node}
	\begin{flushleft}
		
		Sometimes, we encounter a node $m$ that has
		already been discovered.
		
		Let $g(m)$ be its current tentative cost. We
		compare it with the cost of reaching $m$ through
		the node $k$ currently being explored:
		
		\[
		g_{\mathrm{new}}(m)=g(k)+c(k,m).
		\]
		
		There are two possibilities:
		
		\begin{itemize}
			\item If $g(m)\leq g_{\mathrm{new}}(m)$,
			keep the current cost and predecessor.
			\item If $g(m)>g_{\mathrm{new}}(m)$,
			update the cost and set $k$ as the new
			predecessor of $m$:
			\[
			g(m)\leftarrow g_{\mathrm{new}}(m).
			\]
		\end{itemize}
		
		This operation is called \textbf{relaxation}
		of an edge.
		
	\end{flushleft}
\end{frame}

%------------------------------------------------
\begin{frame}{Dijkstra's Algorithm: Selecting the Next Node}
	\begin{flushleft}
		
		The key step is selecting the next node to explore.
		
		We select the node $k$ with the smallest tentative
		cost among all nodes in the open set:
		
		\[
		k=\arg\min_{m\in\mathrm{Open}} g(m).
		\]
		
		We then explore $k$ and move it to the closed set.
		
		\begin{itemize}
			\item If $k$ is the goal node, we terminate.
			\item Otherwise, we continue selecting and
			exploring nodes.
		\end{itemize}
		
		Because all edge costs are nonnegative, the
		cost of a selected node is its shortest-path cost.
		
	\end{flushleft}
\end{frame}



%------------------------------------------------
\begin{frame}{Dijkstra's Algorithm: Recovering the Path}
	\begin{flushleft}
		
		Once the algorithm terminates, we can recover
		the shortest path to the goal. Starting from the goal node, we repeatedly find
		its predecessor and add it to the path.
		
		\medskip
		
		Equivalently, for the current node $m$, we can
		select the predecessor $k$ satisfying
		\[
		k=\arg\min_{k:\,m\in\mathcal{N}(k)}
		\bigl(g(k)+c(k,m)\bigr).
		\]
		
		We continue until we reach the starting node $s$.
		Then we reverse the sequence of nodes to obtain
		the path from $s$ to the goal.
		
	\end{flushleft}
\end{frame}

%------------------------------------------------
\begin{frame}{Dijkstra's Algorithm: One Goal or All Goals?}
	\begin{flushleft}
		
		Dijkstra's algorithm uses the goal node only
		to determine when to terminate.
		
		
		\medskip
		
		If we let the algorithm run until the open set
		becomes empty, it computes the shortest-path
		costs from the starting node to every reachable
		node in the graph.
		
		\begin{itemize}
			\item We can then choose any reachable node
			as the goal and recover its shortest path.
			\item However, this may be unnecessary work
			if we only need a path to one known goal.
		\end{itemize}
		
		We can use knowledge of the goal to guide the
		search and explore fewer nodes.
		
	\end{flushleft}
\end{frame}

%------------------------------------------------
\begin{frame}{A$^\ast$ Algorithm}
	\begin{flushleft}
		
		The A$^\ast$ algorithm is similar to Dijkstra's
		algorithm, but it uses information about the goal
		to guide the selection of nodes for exploration.
		
		\medskip
		While Dijkstra selects the node with the smallest
		tentative cost $g(k)$, A$^\ast$ selects the node
		with the smallest estimated total cost:
		
		\[
		f(k)=g(k)+h(k).
		\]
		
		Here:
		\begin{itemize}
			\item $g(k)$ is the cost of the best path
			found so far from the start to node $k$.
			\item $h(k)$ estimates the remaining cost
			from node $k$ to the goal.
		\end{itemize}
		
		An \textbf{admissible heuristic} never overestimates
		the true minimum remaining cost.
		
	\end{flushleft}
\end{frame}




%------------------------------------------------
\begin{frame}{Value Iteration}
	\begin{flushleft}
		
		Dijkstra's algorithm starts from the initial node
		and propagates information about the cost of
		reaching other nodes.
		
		
		\medskip
		
		Alternatively, we can start from the goal and
		propagate information about the cost of reaching
		the goal from every other node.
		
		\bigskip
		This quantity is called the \textbf{cost-to-go},
		denoted by $V(k)$ for each node $k$. This is the key idea behind \textbf{value iteration}.
		
	\end{flushleft}
\end{frame}

%------------------------------------------------
\begin{frame}{Value Iteration: Initialization}
	\begin{flushleft}
		
		We initialize every node $k$ with a cost-to-go
		estimate $V(k)$.
		
		\begin{itemize}
			\item For the goal node $g$, the cost-to-go is
			zero:
			\[
			V(g)=0.
			\]
			
			\item For all other nodes, the initial estimate
			is infinity:
			\[
			V(k)=+\infty,\qquad k\ne g.
			\]
		\end{itemize}
		
		We then iteratively update these estimates using
		the costs of moving to adjacent nodes.
		
	\end{flushleft}
\end{frame}

%------------------------------------------------
\begin{frame}{Value Iteration: Updates}
	\begin{flushleft}
		
		We update the cost-to-go estimates using the
		\textbf{Bellman equation}:
		
		\[
		V(k)=\min_{m\in\mathcal{N}(k)}
		\bigl(c(k,m)+V(m)\bigr).
		\]
		
		This equation states that the minimum cost to
		the goal is obtained by choosing the best
		successor node and adding the cost of moving
		to that node.
		
		\medskip
		In each iteration, we update the estimates for
		all nodes using the current estimates of their
		neighbors.
		
		\medskip
		We repeat this process until the estimates
		converge.
		
	\end{flushleft}
\end{frame}

%------------------------------------------------
\begin{frame}{Value Iteration: Recovering the Path}
	\begin{flushleft}
		
		Once the algorithm converges, we can recover
		the shortest path from any starting node $s$.
		
		\medskip
		At each node $k$, we select the successor $m$
		that minimizes the sum of the edge cost and
		the cost-to-go:
		
		\[
		m=\arg\min_{m\in\mathcal{N}(k)}
		\bigl(c(k,m)+V(m)\bigr).
		\]
		
		We move to the selected successor and repeat
		until we reach the goal node $g$.
		
	\end{flushleft}
\end{frame}




\begin{frame}{Literature}
	\begin{itemize}
		
		\item Erik Demaine, Ronald Rivest, Srini Devadas -- MIT:
		\bref{https://ocw.mit.edu/courses/6-006-introduction-to-algorithms-spring-2008/pages/lecture-notes/}
		{Introduction to Algorithms: Lecture Notes, lectures 15--18 (Shortest Paths and Dijkstra).}
		
		\item Jason Ku, Justin Solomon -- MIT:
		\bref{https://ocw.mit.edu/courses/6-006-introduction-to-algorithms-spring-2020/resources/mit6_006s20_lec13/}
		{Introduction to Algorithms: Lecture 13, Dijkstra's Algorithm.}
		
		\item MIT OpenCourseWare:
		\bref{https://ocw.mit.edu/courses/6-231-dynamic-programming-and-stochastic-control-fall-2015/resources/mit6_231f15_lec10/}
		{Dynamic Programming and Stochastic Control: Lecture 10, Bellman's Equation and Value Iteration.}
		
	\end{itemize}
\end{frame}




\myqrframe




\end{document}
